\chapter{Обозначения}
\label{sec:notation}

Типы нейронов обозначены первыми четырьмя буквами главного медиатора: \nt{ГЛУТ} --- глутамат, \nt{ГАМК} --- гамма-аминомасляная кислота, \nt{ГЛИЦ} --- глицин, \nt{АЦЕТ} --- ацетилхолин, \nt{ДОФА} --- дофамин, \nt{СЕРО} --- серотонин, \nt{ГИСТ} --- гистамин, \nt{НОРА} --- норадреналин, \nt{АДРЕ} --- адреналин, \nt{АСПА} --- аспартат (таблица~\ref{tab:types}).

Букв меньше, чем величин, и в разных главах одна буква иногда означает разное. В таблице ниже у каждого значения своя строка; в правом столбце --- формула, где оно введено.

\begin{center}
\footnotesize
\setlength{\tabcolsep}{4pt}
\renewcommand{\arraystretch}{1.12}
\begin{longtable}{>{\raggedright}p{2.0cm}>{\raggedright}p{9.6cm}>{\raggedright\arraybackslash}p{2.4cm}}
\toprule
Символ & Значение & Где введено \\
\midrule
\endfirsthead
\toprule
Символ & Значение & Где введено \\
\midrule
\endhead
\bottomrule
\endfoot
$a$ & крутизна линейной передаточной характеристики, $f(u)=au$ & \eqref{eq:feedback} \\
$a$ & уровень \nt{АЦЕТ}: $0$ --- чтение, $1$ --- запись & \eqref{eq:acetnet} \\
$a_i$ & усталость группы $i$ в генераторе ритма & \eqref{eq:halfcentre} \\
$A(r)$, $A_0$ & ответ синапса на импульс при частоте $r$; у отдохнувшего синапса & \eqref{eq:depression} \\
$A_1$ & ответ получателя на одну порцию медиатора & \eqref{eq:quantal} \\
$b_i$ & собственный приток ритмоводителя & \eqref{eq:seroneuron} \\
$b$ & как быстро работа утомляет группу в генераторе ритма & \eqref{eq:halfcentre} \\
$c$ & концентрация медиатора в объёме & \eqref{eq:seroneuron} \\
$c$ & коэффициент корреляции шумов соседних нейронов & \eqref{eq:correlated} \\
$C$ & цена усилия: действие за время $\tau_a$ стоит $C/\tau_a$ & \eqref{eq:vigor} \\
$d'$ & различимость двух входов: разница, делённая на шум & \eqref{eq:gainsnr} \\
$d$, $d_j$ & задержка на пути сигнала & \eqref{eq:glutneuron}, \eqref{eq:vstar} \\
$d$ & задержка петли обратной связи с мышцей & \eqref{eq:delaystable} \\
$D$ & коэффициент диффузии & \eqref{eq:diffusionlength} \\
$D$ & знаменатель в частотах \nt{ГЛУТ}--\nt{ГАМК}-сети & \eqref{eq:isn} \\
$D$ & время ожидания отложенной награды & \eqref{eq:patience} \\
$e$, $e_{ij}$ & след в синапсе & \eqref{eq:threefactor} \\
$E_{\text{имп}}$ & энергия одного импульса вместе с работой синапсов & \eqref{eq:energy} \\
$E$ & частота \nt{ГЛУТ}-группы & \eqref{eq:ei} \\
$E_E$ & равновесный потенциал возбуждающего синапса & \eqref{eq:shunt} \\
$f$, $F$ & передаточная характеристика: частота как функция входа & \eqref{eq:ratenet}, \eqref{eq:gain} \\
$F(t)$, $F_1$ & сила мышцы; сила одного подёргивания & \eqref{eq:forcesum} \\
$F_1$, $F_2$ & силы двух мышц, тянущих сустав в разные стороны & \eqref{eq:antagonists} \\
$g$ & чувствительность рецептора к концентрации & \eqref{eq:seroneuron} \\
$g_L$, $g_E$, $g_I$ & проводимости утечки, возбуждения и торможения & \eqref{eq:shunt} \\
$g$ & коэффициент усиления, задаваемый \nt{НОРА} & \eqref{eq:gain}, \eqref{eq:machine} \\
$g$ & общее торможение от \nt{ГАМК} в сети памяти & \eqref{eq:acetnet} \\
$g$ & усиление ворот боли, задаваемое сверху & \eqref{eq:gate} \\
$G$ & усиление петли обратной связи & \eqref{eq:feedback} \\
$G$ & скорость исправления в петле с задержкой & \eqref{eq:delaystable} \\
$h$, $h_I$ & внешний вход \nt{ГЛУТ}-группы и \nt{ГАМК}-группы & \eqref{eq:ei}, \eqref{eq:isn} \\
$h(t)$ & одиночное подёргивание мышцы & \eqref{eq:twitch} \\
$I$, $I_i$ & внешний приток к нейрону или группе & \eqref{eq:ratenet}, \eqref{eq:wta} \\
$I$ & частота \nt{ГАМК}-группы & \eqref{eq:ei} \\
$I$ & яркость или интенсивность раздражителя & \eqref{eq:photonnoise}, \eqref{eq:nakarushton} \\
$k$ & число импульсов в пачке & \eqref{eq:burst} \\
$k$ & во сколько раз прибавка должна превышать шум & \eqref{eq:photonnoise} \\
$k_s$ & сила, с которой боль накачивает сенсибилизацию & \eqref{eq:sensitization} \\
$K$ & число признаков результата & \eqref{eq:vectortd} \\
$K$ & жёсткость сустава & \eqref{eq:antagonists} \\
$L$ & длина линии задержек & \eqref{eq:itd} \\
$L$ & длина провода от датчика боли & \eqref{eq:paintime} \\
$M$ & мера причинной связи предвестника и награды & \eqref{eq:retro} \\
$M_k$ & ожидание признака $k$ & \eqref{eq:vectortd} \\
$M_{ik}$ & сила торможения от \nt{ГАМК}-нейрона $k$ & \eqref{eq:machine} \\
$M$ & вращающий момент сустава & \eqref{eq:antagonists} \\
$n$, $\bar n$ & число импульсов за окно; его среднее & \eqref{eq:rateerror} \\
$N$ & число нейронов & \eqref{eq:correlated}, \eqref{eq:energy} \\
$N_s$ & число мест выброса между двумя нейронами & \eqref{eq:quantal} \\
$p$ & вероятность выброса медиатора на импульс & \eqref{eq:burst} \\
$p$ & доля активных нейронов в сети памяти & \eqref{eq:acetnet} \\
$P$ & мощность, затрачиваемая на сигналы & \eqref{eq:energy} \\
$P$ & неопределённость хранимой оценки & \eqref{eq:kalman} \\
$q$ & скорость, с которой меняется мир & \eqref{eq:kalman} \\
$q_k$ & частота \nt{ГАМК}-нейрона $k$ & \eqref{eq:machine} \\
$q_0$ & фоновая активность тормозящего посредника & \eqref{eq:gate} \\
$q$ & частота тормозящего посредника в воротах боли & \eqref{eq:gate} \\
$q$ & во сколько раз следующая двигательная единица сильнее предыдущей & \eqref{eq:sizeprinciple} \\
$r$, $r_i$ & частота импульсов нейрона & \eqref{eq:ratenet} \\
$r$, $r_t$ & награда (поток награды) & \eqref{eq:rpe}, \eqref{eq:ctd} \\
$R$, $\bar R$ & полученная награда; её ожидание или среднее в единицу времени & \eqref{eq:perturbation}, \eqref{eq:vigor} \\
$R$ & дисперсия шума на входе фильтра & \eqref{eq:kalman} \\
$R$, $r_0$ & отложенная и немедленная награды & \eqref{eq:patience} \\
$R_{\max}$ & предельная скорость передачи, бит/с & \eqref{eq:spikeentropy} \\
$r_{\max}$ & наибольшая частота импульсов & \eqref{eq:gain}, \eqref{eq:nakarushton} \\
$s_j$ & знак выходов нейрона $j$ & \eqref{eq:dalesign} \\
$s_i$ & знак рецептора получателя & \eqref{eq:seroneuron} \\
$s$, $\hat s$ & состояние мира; его оценка & \eqref{eq:rpe}, \eqref{eq:kalman} \\
$t_1$ & момент первого импульса & \eqref{eq:latency} \\
$T$ & окно счёта импульсов; время накопления фотонов & \eqref{eq:rateerror}, \eqref{eq:photonnoise} \\
$T_c$ & время нарастания подёргивания мышцы & \eqref{eq:twitch} \\
$u$ & суммарный вход нейрона & \eqref{eq:gain} \\
$U$ & уровень постоянного входа & \eqref{eq:latency} \\
$U$ & доля запаса медиатора, выбрасываемая на импульс & \eqref{eq:depression} \\
$v$ & скорость сигнала в проводе; скорость движения пятна & \eqref{eq:itd}, \eqref{eq:vstar} \\
$V$ & напряжение на мембране & \eqref{eq:glutneuron}, \eqref{eq:shunt} \\
$V$ & ожидание будущей награды & \eqref{eq:rpe}, \eqref{eq:ctd} \\
$w$, $w_{ij}$, $W_{ij}$ & вес связи & \eqref{eq:glutneuron}, \eqref{eq:ratenet}, \eqref{eq:acetnet} \\
$w_{q\mu}$, $w_{q\nu}$, $w_{z\mu}$ & веса связей в воротах боли & \eqref{eq:gate} \\
$x$, $y$ & логические входы и выход & \eqref{eq:dualrail}, \eqref{eq:veto} \\
$x$, $y$ & активность отправителя и получателя & \eqref{eq:hebb} \\
$x$ & место детектора на линии задержек & \eqref{eq:itd} \\
$x_i$ & вход от органов чувств & \eqref{eq:acetnet} \\
$z$ & частота выхода ворот боли & \eqref{eq:gate} \\
$y$ & зашумлённый вход фильтра & \eqref{eq:kalman} \\
\midrule
$\alpha_i^\pm$ & вес положительных и отрицательных ошибок & \eqref{eq:asymmetric} \\
$\beta$ & сила взаимного торможения & \eqref{eq:wta} \\
$\beta$ & сила торможения выхода ворот боли & \eqref{eq:gate} \\
$\gamma$ & коэффициент дисконтирования & \eqref{eq:rpe} \\
$\delta(\cdot)$ & дельта-функция: мгновенный скачок & \eqref{eq:glutneuron} \\
$\delta$, $\delta_t$ & ошибка предсказания награды & \eqref{eq:rpe}, \eqref{eq:ctd} \\
$\delta t$ & разница времён прихода звука в два уха & \eqref{eq:itd} \\
$\Delta$, $\Delta_{\max}$ & интервал между импульсами; окно совпадения & \eqref{eq:window} \\
$\Delta t$ & точность, с которой получатель различает время & \eqref{eq:spikeentropy} \\
$\Delta u$ & разница входов двух групп & \eqref{eq:gainsnr} \\
$\Delta x$ & расстояние между соседними датчиками & \eqref{eq:vstar} \\
$\varepsilon$ & относительная разница входов & \eqref{eq:wtagain} \\
$\eta$ & скорость обучения & \eqref{eq:plasticity} \\
$\eta$ & шум в сети после усилителя & \eqref{eq:softmax} \\
$\mu$ & вход прикосновения к воротам боли & \eqref{eq:gate} \\
$\nu$ & вход боли & \eqref{eq:gate} \\
$\theta$ & порог срабатывания & \eqref{eq:window} \\
$\kappa$ & количество медиатора на один импульс & \eqref{eq:seroneuron} \\
$\kappa_i$ & отношение весов положительных и отрицательных ошибок & \eqref{eq:expectile} \\
$\kappa_m$ & связь силы мышцы с её жёсткостью & \eqref{eq:antagonists} \\
$\ell$ & расстояние, на которое расплывается медиатор & \eqref{eq:diffusionlength} \\
$\ell$ & плечо, на котором мышца тянет сустав & \eqref{eq:antagonists} \\
$\xi$ & случайное дрожание выхода нейрона & \eqref{eq:perturbation} \\
$\sigma$ & разброс, масштаб шума & \eqref{eq:rateerror}, \eqref{eq:softmax} \\
$\sigma$ & яркость, при которой ответ равен половине наибольшего & \eqref{eq:nakarushton} \\
$\sigma_{\text{до}}$, $\sigma_{\text{после}}$ & шум до усилителя и после него & \eqref{eq:gainsnr} \\
$\tau$ & постоянная времени мембраны & \eqref{eq:glutneuron} \\
$\tau_c$ & время жизни медиатора в объёме & \eqref{eq:seroneuron} \\
$\tau_E$, $\tau_I$ & постоянные времени \nt{ГЛУТ}- и \nt{ГАМК}-группы & \eqref{eq:ei} \\
$\tau_e$ & время жизни следа & \eqref{eq:threefactor} \\
$\tau_s$ & время возврата чувствительности после повреждения & \eqref{eq:sensitization} \\
$\tau_{\text{вос}}$ & время восстановления запаса медиатора & \eqref{eq:depression} \\
$\tau_\gamma$ & горизонт планирования & \eqref{eq:ctd} \\
$\tau_v$ & скорость, с которой уточняется ожидание & \eqref{eq:ramp} \\
$\tau_a$ & время усталости в генераторе ритма & \eqref{eq:halfcentre} \\
$\tau_a^*$ & лучшее время выполнения действия & \eqref{eq:vigor} \\
$\Phi$ & правая часть правила обучения & \eqref{eq:plasticity} \\
$\Phi$ & поток фотонов & \eqref{eq:photonnoise} \\
$\varphi_k$ & признак $k$ полученного результата & \eqref{eq:vectortd} \\
$\boldsymbol\psi$ & настройка сети в модели DishBrain & \eqref{eq:dishdensity} \\
$\rho(\boldsymbol\psi)$ & доля времени, проведённого в настройке $\boldsymbol\psi$ & \eqref{eq:dishdensity} \\
$P_{\text{пр}}$ & вероятность промаха & \eqref{eq:dishdensity} \\
$n_\uparrow$, $n_\downarrow$ & счёт импульсов областей «вверх» и «вниз» за окно & \eqref{eq:dishreadout} \\
$c$ & коэффициент, с которым разность счётов двигает ракетку & \eqref{eq:dishreadout} \\
\end{longtable}
\end{center}
